\documentclass[10pt,a4paper]{article}
\usepackage{cv-theme}
\begin{document}
{\LARGE\bfseries Jacob M. Graving\par}\vspace{5pt}
\href{mailto:jgraving@gmail.com}{jgraving@gmail.com} | \href{https://jakegraving.com/}{jakegraving.com} | \href{https://github.com/jgraving}{github.com/jgraving}\par
Max Planck Institute of Animal Behavior · University of Konstanz, Department of Biology\par
Universitätsstr. 10, Konstanz, Germany 78464\par
\section*{Summary}
Experienced research scientist and machine learning engineer with a strong foundation in data-driven modeling, statistical thinking, and problem-solving at scale. Comfortable moving between research and engineering, theory and practice. Passionate about building systems that turn complex data into actionable insight.\par
\section*{Skills}
\begin{itemize}
\item Design and implementation of probabilistic models and Bayesian inference pipelines for complex, uncertain data
\item Development of custom deep learning models for behavior, vision, and sequences using PyTorch, and JAX
\item Building scalable, reproducible pipelines for multimodal and high-dimensional data (ETL, analysis, visualization)
\item Creating interpretable ML tools that support scientific discovery and real-world decision-making
\item Leading and contributing to open-source software projects in Python; experienced with PyMC, NumPyro, Stan, scikit-learn, and more
\item Strong communicator and collaborator across disciplines - biology, neuroscience, engineering, and data science
\end{itemize}
\section*{Experience}
\subsection*{2020-present | Research Scientist}
Advanced Research Technology Unit, Max Planck Institute of Animal Behavior\par
\begin{itemize}
\item Lead independent research at the interface of machine learning and behavioral science, developing general-purpose tools to measure and model animal behavior in both lab and field environments.
\item Design and maintain scalable ETL pipelines for multimodal behavioral datasets, ensuring data integrity and reproducibility across experimental systems.
\item Collaborate with interdisciplinary teams - including software engineers, neuroscientists, and field biologists - to build AI-driven tools for behavior analysis and experimental optimization.
\item Apply Bayesian inference, causal modeling, and uncertainty quantification to large-scale time series and experimental data to generate interpretable, data-driven insights.
\item Provide technical guidance and mentorship across collaborative projects, promoting reproducible science and statistical rigor.
\end{itemize}
\section*{Education}
\subsection*{2021 | Dr.rer.nat., Biology (0.0 ‘summa cum laude’)}
Max Planck Institute of Animal Behavior \& University of Konstanz  -  Germany\par
International Max Planck Research School (IMPRS)\par
Thesis: \href{http://nbn-resolving.de/urn:nbn:de:bsz:352-2-dgcbudqch6ix8}{Deep Learning and Computer Vision Methods for Measuring and Modeling Animal Behavior}\par
\subsection*{2015 | M.Sc., Biology}
Bowling Green State University  -  USA\par
Focus: Animal Behavior, Neuroscience\par
\subsection*{2013 | B.Sc., Biology}
Bowling Green State University  -  USA\par
\newpage
\section*{Recent \& Selected Publications}
\subsection*{In Prep}
\textbf{Graving, J.M.} and Foster, J.J. (in prep). Unwrapping Circular Statistics: Bayesian Linear Models for Circular Data.\par
\begin{itemize}
\item Introduces a novel Bayesian generalized linear model framework for circular data by incorporating the von Mises distribution.
\item Jointly models mean direction and variance to improve interpretability of circular outcomes.
\item Provides practical guidelines and implementations using Bayesian toolkits such as PyMC and Stan.
\end{itemize}
\textbf{Graving, J.M.}, Heins, C., Couzin, I.D. Revealing the Structure of Time-Series Data with Context Attraction-Repulsion Embeddings.\par
\begin{itemize}
\item Developed a generalized contrastive learning framework for dimensionality reduction and visualization of time-series data using transformer-based sequence modeling.
\item Demonstrated recovery of latent factors and interpretable structure from high-dimensional behavioral and synthetic datasets.
\item Provides a general tool for exploratory analysis and hypothesis generation in time-series modeling.
\end{itemize}
\subsection*{In Review}
Bath, D.E., \textbf{Graving, J.M.}, Walter, T., Sridhar, V.H., Vizcaíno, J.P., Couzin, I.D. Collective detection and processing of distributed information by fish schools. In revision for \textit{Current Biology}.\par
Code: \href{https://github.com/jgraving/bayesian_beta_regression}{github.com/jgraving/bayesian\_beta\_regression}\par
\begin{itemize}
\item Led development of a Bayesian model to analyze collective information processing in fish schools.
\item Modeled 3 billion behavioral observations across 400+ conditions using JAX and NumPyro.
\item Revealed how distributed sensing shapes group-level responses to environmental signals.
\end{itemize}
\subsection*{2025}
Sayin, S., Couzin-Fuchs, E., Petelski, I., Günzel, Y., Salahshour, M., Lee, C.-Y., \textbf{Graving, J.M.}, Li, L., Deussen, O., Sword, G.A., Couzin, I.D. (2025). The behavioral mechanisms governing collective motion in swarming locusts. \textit{Science}, 387(6737), 995-1000. \href{https://doi.org/10.1126/science.adq7832}{https://doi.org/10.1126/science.adq7832}\par
\begin{itemize}
\item Led Bayesian modeling and statistical analysis of VR behavioral data from locusts.
\item Identified and solved a key methodological flaw in prior work - confounding of group coordination with group size.
\item Helped validate the core findings and strengthen the paper’s quantitative rigor.
\end{itemize}
\subsection*{2023}
Koger, B., Deshpande, A., Kerby, J.T., \textbf{Graving, J.M.}, Costelloe, B. R., Couzin, I.D. (2023). Quantifying the movement, behaviour and environmental context of group-living animals using drones and computer vision. \textit{Journal of Animal Ecology}. \href{https://doi.org/10.1111/1365-2656.13904}{doi:10.1111/1365-2656.13904} Code: \href{https://github.com/benkoger/overhead-video-worked-examples}{github.com/benkoger/overhead-video-worked-examples}\par
\begin{itemize}
\item Co-developed a computer vision pipeline combining drone footage and deep learning to extract 3D behavioral and environmental data from animal groups in natural settings.
\item Integrated spatial context reconstruction with fine-scale movement tracking to support large-scale behavioral ecology studies.
\end{itemize}
\newpage
\section*{Recent \& Selected Publications (continued)}
\subsection*{2020}
Li, L., Nagy, M., \textbf{Graving, J.M.}, Bak-Coleman, J., Guangming X., Couzin, I.D. (2020). Vortex phase matching as a strategy for schooling in robots and in fish. \textit{Nature Communications}, 11, 5408. \href{https://doi.org/10.1038/s41467-020-19086-0}{doi:10.1038/s41467-020-19086-0}\par
\begin{itemize}
\item Analyzed high-resolution posture data of schooling fish to test energy-optimization strategies predicted by robotic models.
\item Applied deep learning-based pose estimation and mechanistic modeling to quantify fluid-mediated interactions and coordinated motion patterns.
\end{itemize}
\subsection*{2019}
\textbf{Graving, J.M.}, Chae, D., Naik, H., Li, L., Koger, B., Costelloe, B.R., Couzin, I.D. (2019). DeepPoseKit, a software toolkit for fast and robust animal pose estimation using deep learning. \textit{eLife}, 8. \href{https://doi.org/10.7554/elife.47994}{doi:10.7554/elife.47994}\par
Press: \href{https://www.quantamagazine.org/to-decode-the-brain-scientists-automate-the-study-of-behavior-20191210/}{Quanta Magazine}, \href{https://doi.org/10.1038/s41592-019-0678-2}{Nature Methods}, \href{https://doi.org/10.1038/d41586-019-02942-5}{Nature News \& Views}, \href{https://elifesciences.org/digests/47994/machine-learning-animal-poses-to-understand-behavior}{eLife Science Digests}\par
Code: \href{https://github.com/jgraving/deepposekit}{github.com/jgraving/deepposekit}\par
\begin{itemize}
\item Created a general-purpose, few-shot deep learning framework for high-speed, high-accuracy animal pose tracking in Python/TensorFlow.
\item Supervised development of key components including a custom annotation GUI and augmentation tools for low-data regimes.
\end{itemize}
\section*{Teaching}
\subsection*{2023 | Konstanz School of Collective Behavior}
University of Konstanz\par
\begin{itemize}
\item Designed and led a workshop on “Probabilistic Machine Learning” for 30+ international PhD students, covering Bayesian inference, causal inference, and information theory in practice. \href{https://www.exc.uni-konstanz.de/kscb/}{https://www.exc.uni-konstanz.de/kscb/}
\end{itemize}
\subsection*{2022 | Deep Learning Workshop}
Max Planck Institute of Animal Behavior\par
\begin{itemize}
\item Led an applied workshop introducing PyTorch and deep learning fundamentals to 15 researchers.
\end{itemize}
\subsection*{2019 | ASAB 2019 Summer Conference}
University of Konstanz\par
\begin{itemize}
\item Gave an invited seminar on “Machine Learning in the Behavioral Sciences” to an audience of 400+.
\item Co-organized a workshop on behavioral quantification and modeling using ML methods (27 participants).
\end{itemize}
\subsection*{2016-2020 | University of Konstanz, Department of Biology}
\begin{itemize}
\item Co-developed and taught an intensive research course on collective behavior for Master’s students.
\item Supervised five Master’s thesis projects across behavior, computation, and data science.
\item Delivered lectures and practicals on computer vision for behavior analysis, data processing, and Python programming.
\end{itemize}
\end{document}
